%%%%%%%%%%%%%%%%%%%%%%%%%%%%%%%%%%%%%%%%%
% Medium Length Professional CV
% LaTeX Template
% Version 3.0 (December 17, 2022)
%
% This template originates from:
% https://www.LaTeXTemplates.com
%
% Author:
% Vel (vel@latextemplates.com)
%
% Original author:
% Trey Hunner (http://www.treyhunner.com/)
%
% License:
% CC BY-NC-SA 4.0 (https://creativecommons.org/licenses/by-nc-sa/4.0/)
%
%%%%%%%%%%%%%%%%%%%%%%%%%%%%%%%%%%%%%%%%%

%----------------------------------------------------------------------------------------
%	PACKAGES AND OTHER DOCUMENT CONFIGURATIONS
%----------------------------------------------------------------------------------------

\documentclass[
	%a4paper, % Uncomment for A4 paper size (default is US letter)
	11pt, % Default font size, can use 10pt, 11pt or 12pt
]{resume} % Use the resume class

\usepackage{ebgaramond} % Use the EB Garamond font
\usepackage{hyperref}

\hypersetup{
    colorlinks=true,
    linkcolor=blue,
    filecolor=magenta,      
    urlcolor=blue,
    pdftitle={Overleaf Example},
    pdfpagemode=FullScreen,
    }

%------------------------------------------------

\name{Tushar Mopuri} % Your name to appear at the top

% You can use the \address command up to 3 times for 3 different addresses or pieces of contact information
% Any new lines (\\) you use in the \address commands will be converted to symbols, so each address will appear as a single line.

\address{Fourth-Year PhD Student \\ \href{https://tmopuri.com/}{Website}}
\address{tmopuri@upenn.edu} % Contact information

%----------------------------------------------------------------------------------------

\begin{document}

%----------------------------------------------------------------------------------------
%	EDUCATION SECTION
%----------------------------------------------------------------------------------------

\begin{rSection}{Education}
	
	\textbf{University of Pennsylvania} \hfill \textit{August 2023 - Present} \\ 
	PhD in Computer and Information Science \\
	Advisor: Pratyush Mishra

        \textbf{Indian Institute of Science} \hfill \textit{August 2018 - April 2022} \\ 
	BS (Research) \\
        Major: Mathematics
	
\end{rSection}


%----------------------------------------------------------------------------------------
%	WORK EXPERIENCE SECTION
%----------------------------------------------------------------------------------------

\begin{rSection}{Preprints}

        Authors are listed in alphabetical order (by last name), unless otherwise noted.

        \textbf{Vector Commit-and-Reduce Schemes} \\ 
	Anubhav Baweja, Pratyush Mishra, Tushar Mopuri, Bharath Namboothiry \\
        (\href{https://eprint.iacr.org/2026/2790}{ePrint})

\end{rSection}


%----------------------------------------------------------------------------------------
%	WORK EXPERIENCE SECTION
%----------------------------------------------------------------------------------------

\begin{rSection}{Conference Publications}

        \textbf{Field-Agnostic SNARKs with Small Proofs via Encode-Repeat-Accumulate (ERA) Codes} \\ 
	Anubhav Baweja, Giacomo Fenzi, Pratyush Mishra, Tushar Mopuri \\
        $\mathbf{ASIACRYPT} \; \mathbf{2026}$ \\
        Also presented at ZKProof 8 \\
        (\href{https://eprint.iacr.org/2026/864}{ePrint})

        \textbf{FICS and FACS: Fast IOPPs and Accumulation via Code-Switching} \\ 
	Anubhav Baweja, Pratyush Mishra, Tushar Mopuri, Matan Shtepel \\
        $\mathbf{CRYPTO} \; \mathbf{2026}$ \\
        Also presented at ZKProof 8 (merged) \\
        (\href{https://eprint.iacr.org/2025/737}{ePrint})

        \textbf{Query-Optimal IOPPs for Linear-Time Encodable Codes} \\ 
	Anubhav Baweja, Pratyush Mishra, Tushar Mopuri, Matan Shtepel \\
        $\mathbf{EUROCRYPT} \; \mathbf{2026}$ \\
        Also presented at ZKProof 8 (merged) \\
        (\href{https://eprint.iacr.org/2025/1588}{ePrint})

        \textbf{Scribe: Low-memory SNARKs via Read-Write Streaming} \\ 
	Anubhav Baweja, Pratyush Mishra, Tushar Mopuri, Karan Newatia, Steve Wang \\
        $\mathbf{USENIX} \; \mathbf{Security} \; \mathbf{2026}$ \\
        \textbf{Runner-Up for Distinguished Paper Award} \\
        Also presented at ZKProof 8 and zkSummit 11 \\
        (\href{https://eprint.iacr.org/2024/1970}{ePrint})

        \textbf{Time-Space Trade-Offs for Sumcheck} \\ 
	Anubhav Baweja, Alessandro Chiesa, Elisabetta Fedele, Giacomo Fenzi, Pratyush Mishra, Tushar Mopuri, Andrew Zitek-Estrada \\
        $\mathbf{TCC} \; \mathbf{2025}$ \\
        Also presented at ZKProof 8 and zkSummit 14 \\
        (\href{https://eprint.iacr.org/2025/1473}{ePrint})

        \textbf{DewTwo: a transparent PCS with quasi-linear prover, logarithmic verifier and 4.5KB proofs from falsifiable assumptions} \\ 
	Benedikt Bünz, Tushar Mopuri, Alireza Shirzad, Sriram Sridhar \\
        $\mathbf{CRYPTO} \; \mathbf{2025}$ \\
        (\href{https://eprint.iacr.org/2025/129}{ePrint})

        \textbf{Hekaton: Horizontally-Scalable zkSNARKs via Proof Aggregation} \\ 
	Michael Rosenberg, Tushar Mopuri, Hossein Hafezi, Ian Miers, Pratyush Mishra (not alphabetical) \\
        $\mathbf{ACM} \; \mathbf{CCS} \; \mathbf{2024}$ \\
        (\href{https://eprint.iacr.org/2024/1208}{ePrint})

        \textbf{Dew: Transparent Constant-sized zkSNARKs} \\ 
	Arasu Arun, Chaya Ganesh, Satya Lokam, Tushar Mopuri, Sriram Sridhar \\
        $\mathbf{PKC} \; \mathbf{2023}$ \\
        (\href{https://eprint.iacr.org/2022/419}{ePrint})

\end{rSection}

%----------------------------------------------------------------------------------------
%	WORK EXPERIENCE SECTION
%----------------------------------------------------------------------------------------

\begin{rSection}{Research Visits}

        \textbf{Simons Institute, Berkeley} - Attended cryptography workshops \hfill \textit{Summer 2025} \\ 
        \textbf{New York University} - Worked with Benedikt Bünz \hfill \textit{May 2024 - August 2024} \\ 
        \textbf{Aarhus University} - Worked with Claudio Orlandi, Peter Scholl and Lawrence Roy \hfill \textit{Fall 2022} \\ 
        \textbf{Microsoft Research (India)} - Interned with Satya Lokam \hfill \textit{Spring 2022}

\end{rSection}


%----------------------------------------------------------------------------------------
%	WORK EXPERIENCE SECTION
%----------------------------------------------------------------------------------------

\begin{rSection}{Fellowships}

        \textbf{Jonathan M. Smith Fellowship} \hfill \textit{August 2023 - July 2024} \\ 
	Awarded by the University of Pennsylvania.

        \textbf{Kishore Vaigyanik Protsahan Yojana (KVPY) Fellowship} \hfill \textit{August 2018 - April 2022} \\ 
	Awarded by the Department of Science and Technology (DST), Govt. of India.

\end{rSection}

%----------------------------------------------------------------------------------------
%	TECHNICAL STRENGTHS SECTION
%----------------------------------------------------------------------------------------

\begin{rSection}{Teaching and Reviewing}

    I was a TA for CIS 5560 (Cryptography) at Penn in Spring 2024 and Spring 2025. Additionally, I was an external reviewer for Eurocrypt 2025, Crypto 2025, TCC 2026, Asiacrypt 2026 and USENIX Security 2027.

\end{rSection}

%----------------------------------------------------------------------------------------
%	EXAMPLE SECTION
%----------------------------------------------------------------------------------------

%\begin{rSection}{Section Name}

	%Section content\ldots

%\end{rSection}

%----------------------------------------------------------------------------------------

\end{document}
